\chapter{Appendices}


\section{The Penalty Approach to Constrained Optimization Problems}

\subsection{An Equality-Constrained Problem}

Consider the constrained maximization problem
\[
\max_x f(x)
\]
subject to
\[
g_m(x)=0,
\qquad
m=1,\dots,M.
\]

Suppose that
\[
x^*
\]
is a solution to this problem.

For
\[
j=1,2,\dots,
\]
define the penalized objective function
\[
F_j(x)
=
f(x)
-
\frac{j}{2}\|g(x)\|^2
-
\alpha\|x-x^*\|^2,
\qquad
\alpha>0.
\]

Here,
\[
g(x)
=
(g_1(x),\dots,g_M(x))^\top,
\]
and
\[
\|g(x)\|^2
=
\sum_{m=1}^M g_m(x)^2.
\]

Define the closed set
\[
S
=
\left\{
x:\|x-x^*\|\le \varepsilon
\right\}.
\]

Let
\[
x^j
\]
be an optimal solution of the penalized unconstrained problem
\[
\max_{x\in S}F_j(x).
\]

An optimal solution exists by the Weierstrass theorem because:

\begin{itemize}
    \item the set \(S\) is compact,
    \item the function \(F_j(x)\) is continuous on \(S\).
\end{itemize}

\paragraph{Claim 1.}

The sequence
\[
\{x^j\}
\]
converges to
\[
x^*.
\]

\paragraph{Proof.}

Since
\[
x^j
\]
is optimal,
\[
F_j(x^j)\ge F_j(x)
\]
for every
\[
x\in S.
\]

In particular,
\[
F_j(x^j)\ge F_j(x^*).
\]

Hence,
\[
f(x^j)
-
\frac{j}{2}\|g(x^j)\|^2
-
\alpha\|x^j-x^*\|^2
\ge
f(x^*).
\]

Rearranging,
\[
\frac{j}{2}\|g(x^j)\|^2
+
\alpha\|x^j-x^*\|^2
\le
f(x^j)-f(x^*).
\]

Because
\[
x^j\in S,
\]
the sequence
\[
\{x^j\}
\]
remains in a compact neighborhood of
\[
x^*.
\]

Therefore,
\[
f(x^j)
\]
is bounded above on \(S\).

Suppose
\[
\|g(x^j)\|^2
\]
does not converge to zero.

Then there exists some
\[
\delta>0
\]
and a subsequence such that
\[
\|g(x^j)\|^2\ge\delta.
\]

In that case,
\[
\frac{j}{2}\|g(x^j)\|^2
\ge
\frac{j}{2}\delta
\to\infty
\]
as
\[
j\to\infty,
\]
contradicting the inequality above.

Therefore,
\[
\lim_{j\to\infty}\|g(x^j)\|^2=0.
\]

Hence every limit point
\[
\bar x
\]
of the sequence
\[
\{x^j\}
\]
satisfies
\[
g(\bar x)=0.
\]

Taking limits in the inequality,
\[
f(\bar x)
-
\alpha\|\bar x-x^*\|^2
\ge
f(x^*).
\]

Since
\[
\bar x
\]
is feasible and
\[
x^*
\]
solves the constrained maximization problem,
\[
f(x^*)\ge f(\bar x).
\]

Combining the two inequalities,
\[
f(\bar x)
-
\alpha\|\bar x-x^*\|^2
\ge
f(x^*)
\ge
f(\bar x).
\]

Thus,
\[
-\alpha\|\bar x-x^*\|^2\ge0.
\]

Because
\[
\alpha>0,
\]
this implies
\[
\|\bar x-x^*\|^2=0.
\]

Therefore,
\[
\bar x=x^*.
\]

Every convergent subsequence has the same limit, and therefore
\[
x^j\to x^*.
\]

Hence,
\[
x^*
\]
is an unconstrained local maximizer of
\[
F_j(x)
\]
for sufficiently large
\[
j.
\]

\hfill\(\blacksquare\)

\subsection{Relation to the Lagrangian}

The Lagrangian of the original constrained optimization problem is
\[
\mathcal L(x,\lambda)
=
f(x)-\lambda^\top g(x),
\]
where
\[
\lambda=(\lambda_1,\dots,\lambda_M)^\top
\]
is the vector of Lagrange multipliers.

The first-order condition is
\[
\nabla f(x^*)
-
\nabla g(x^*)\lambda^*
=
0.
\]

Here,
\[
\nabla g(x^*)
\]
is the
\[
N\times M
\]
Jacobian matrix of the constraints:
\[
\nabla g(x^*)
=
\begin{bmatrix}
\nabla g_1(x^*) &
\cdots &
\nabla g_M(x^*)
\end{bmatrix}.
\]

Premultiplying both sides by
\[
[\nabla g(x^*)^\top\nabla g(x^*)]^{-1}
\nabla g(x^*)^\top,
\]
we obtain
\[
[\nabla g(x^*)^\top\nabla g(x^*)]^{-1}
\nabla g(x^*)^\top
\nabla f(x^*)
=
[\nabla g(x^*)^\top\nabla g(x^*)]^{-1}
\nabla g(x^*)^\top
\nabla g(x^*)
\lambda^*.
\]

Since
\[
[\nabla g(x^*)^\top\nabla g(x^*)]^{-1}
\nabla g(x^*)^\top
\nabla g(x^*)
=
I,
\]
we derive
\[
\lambda^*
=
[\nabla g(x^*)^\top\nabla g(x^*)]^{-1}
\nabla g(x^*)^\top
\nabla f(x^*).
\]

The matrix
\[
\nabla g(x^*)^\top\nabla g(x^*)
\]
is invertible because the constraint qualification is satisfied and
\[
\operatorname{rank}(\nabla g(x^*))=M.
\]

Thus, the Lagrange multiplier vector is the projection of the gradient of the objective function onto the linear space generated by the gradients of the constraints.

The penalty approach approximates the constrained optimization problem by a sequence of unconstrained problems in which constraint violations are penalized increasingly heavily as
\[
j\to\infty.
\]

\paragraph{Proposition 1. Theorem of Lagrange}

Let
\[
x^*
\]
be a local maximizer of
\[
f
\]
subject to
\[
g(x)=0.
\]

Assume that the constraint gradients
\[
\nabla g_1(x^*),
\dots,
\nabla g_M(x^*)
\]
are linearly independent.

Then there exists a unique vector
\[
\lambda^*
\]
such that
\[
\nabla f(x^*)
-
\nabla g(x^*)\lambda^*
=
0.
\]

\paragraph{Proof.}

Since
\[
x^j
\]
is a local maximizer of
\[
F_j(x),
\]
the first-order condition implies
\[
\nabla F_j(x^j)=0
\]
for sufficiently large
\[
j.
\]

Therefore,
\[
\nabla f(x^j)
-
j\nabla g(x^j)g(x^j)
-
\alpha(x^j-x^*)
=
0.
\]

As
\[
j\to\infty
\]
and
\[
x^j\to x^*,
\]
we obtain
\[
\nabla f(x^*)
-
\nabla g(x^*)
\lambda^*
=
0,
\]
where
\[
\lambda^*
=
\lim_{j\to\infty}jg(x^j).
\]

Moreover,
\[
\lambda^*
=
[\nabla g(x^*)^\top\nabla g(x^*)]^{-1}
\nabla g(x^*)^\top
\nabla f(x^*).
\]

Therefore, the Lagrange multiplier vector exists and is unique.

\hfill\(\blacksquare\)

\paragraph{Proposition 2. Second-Order Sufficient Condition}

Assume that
\[
f
\]
and
\[
g
\]
are twice continuously differentiable.

Let
\[
x^*\in\mathbb R^N
\]
and
\[
\lambda^*\in\mathbb R^M
\]
satisfy
\[
\nabla_x\mathcal L(x^*,\lambda^*)=0
\]
and
\[
\nabla_\lambda\mathcal L(x^*,\lambda^*)=0.
\]

Suppose also that
\[
z^\top
\nabla_{xx}^2\mathcal L(x^*,\lambda^*)
z
<
0
\]
for every
\[
z\neq0
\]
such that
\[
\nabla g(x^*)^\top z=0.
\]

Then
\[
x^*
\]
is a local maximizer of
\[
f
\]
subject to
\[
g(x)=0.
\]

\paragraph{Proof.}

Consider the Hessian matrix of the penalized objective:
\[
\nabla^2F_j(x^j)
=
\nabla^2f(x^j)
-
j\nabla g(x^j)\nabla g(x^j)^\top
-
j\sum_{m=1}^M
g_m(x^j)\nabla^2g_m(x^j)
-
\alpha I.
\]

Fix any vector
\[
z
\]
such that
\[
\nabla g(x^*)^\top z=0.
\]

Define
\[
z^j
=
z
-
\nabla g(x^j)
[
\nabla g(x^j)^\top
\nabla g(x^j)
]^{-1}
\nabla g(x^j)^\top z.
\]

Then
\[
\nabla g(x^j)^\top z^j=0
\]
and
\[
z^j\to z
\]
as
\[
x^j\to x^*.
\]

Because
\[
\nabla^2F_j(x^j)
\]
is negative definite,
\[
(z^j)^\top
\nabla^2F_j(x^j)
z^j
<
0.
\]

Hence,
\[
(z^j)^\top
\left[
\nabla^2f(x^j)
-
j\sum_{m=1}^M
g_m(x^j)\nabla^2g_m(x^j)
\right]
z^j
-
\alpha\|z^j\|^2
<
0.
\]

Since
\[
z^j\to z
\]
and
\[
jg_m(x^j)\to\lambda_m^*,
\]
we obtain
\[
z^\top
\left[
\nabla^2f(x^*)
-
\sum_{m=1}^M
\lambda_m^*
\nabla^2g_m(x^*)
\right]
z
-
\alpha\|z\|^2
\le0.
\]

Because
\[
\alpha>0
\]
can be arbitrarily small,
\[
z^\top
\left[
\nabla^2f(x^*)
-
\sum_{m=1}^M
\lambda_m^*
\nabla^2g_m(x^*)
\right]
z
<
0.
\]

Therefore,
\[
z^\top
\nabla_{xx}^2
\mathcal L(x^*,\lambda^*)
z
<
0
\]
for every feasible direction
\[
z\neq0
\]
satisfying
\[
\nabla g(x^*)^\top z=0.
\]

Hence,
\[
x^*
\]
is a local maximizer.

\hfill\(\blacksquare\)

\subsection{An Inequality-Constrained Problem}

\paragraph{Proposition 3. Karush--Kuhn--Tucker Necessary Conditions}

Let
\[
x^*
\]
be a local maximizer of
\[
\max_x f(x)
\]
subject to
\[
g(x)=0,
\qquad
h(x)\le0,
\]
where
\[
f,
\quad
g_m,
\quad
h_k
\]
are differentiable functions.

Assume that
\[
x^*
\]
is regular in the sense that the equality constraint gradients
\[
\nabla g_m(x^*)
\]
and the active inequality constraint gradients
\[
\nabla h_k(x^*)
\]
are linearly independent.

Then there exist unique Lagrange multiplier vectors
\[
\lambda^*
\]
and
\[
\mu^*
\]
such that
\[
\nabla_x
\mathcal L(x^*,\lambda^*,\mu^*)
=
0,
\]
\[
\mu_k^*h_k(x^*)=0,
\]
and
\[
\mu_k^*\ge0,
\qquad
k=1,\dots,K.
\]


\paragraph{Proof.}

Consider the penalized problem
\[
\max_{x\in S} F_j(x),
\]
where
\[
F_j(x)
=
f(x)
-
\frac{j}{2}\|g(x)\|^2
-
\frac{j}{2}
\sum_{k=1}^K
\left(h_k^+(x)\right)^2
-
\alpha\|x-x^*\|^2,
\]
with
\[
h_k^+(x)=\max\{0,h_k(x)\}.
\]

Note that
\[
\left(h_k^+(x)\right)^2
\]
is continuously differentiable, and its gradient is
\[
\nabla\left(h_k^+(x)\right)^2
=
2h_k^+(x)\nabla h_k(x).
\]

If
\[
x^j
\]
maximizes
\[
F_j(x)
\]
over
\[
S,
\]
then an argument similar to the equality-constrained case shows that
\[
x^j\to x^*.
\]

The Lagrange multipliers are given by
\[
\lambda_m^*
=
\lim_{j\to\infty}
jg_m(x^j),
\qquad
m=1,\dots,M,
\]
and
\[
\mu_k^*
=
\lim_{j\to\infty}
jh_k^+(x^j),
\qquad
k=1,\dots,K.
\]

Since
\[
h_k^+(x^j)\ge0,
\]
we obtain
\[
\mu_k^*\ge0
\]
for all
\[
k=1,\dots,K.
\]

Also, if
\[
h_k(x^*)<0,
\]
then for sufficiently large \(j\),
\[
h_k^+(x^j)=0,
\]
and hence
\[
\mu_k^*=0.
\]

Therefore,
\[
\mu_k^*h_k(x^*)=0.
\]

\hfill\(\blacksquare\)

\paragraph{Proposition 4. Second-Order Sufficiency Conditions}

Assume that
\[
f,
\quad
g,
\quad
h
\]
are twice continuously differentiable.

Let
\[
x^*\in\mathbb{R}^N,
\qquad
\lambda^*\in\mathbb{R}^M,
\qquad
\mu^*\in\mathbb{R}_+^K
\]
satisfy
\[
\nabla_x\mathcal L(x^*,\lambda^*,\mu^*)=0,
\]
\[
g(x^*)=0,
\]
\[
h(x^*)\le0,
\]
\[
\mu_k^*\ge0,
\]
and
\[
\mu_k^*h_k(x^*)=0.
\]

Suppose that
\[
z^\top
\nabla_{xx}^2
\mathcal L(x^*,\lambda^*,\mu^*)
z
<
0
\]
for all
\[
z\ne0
\]
such that
\[
\nabla g_m(x^*)^\top z=0,
\qquad
m=1,\dots,M,
\]
and
\[
\nabla h_k(x^*)^\top z=0
\]
for every active inequality constraint
\[
k
\]
such that
\[
h_k(x^*)=0.
\]

Assume also that
\[
\mu_k^*>0
\]
for all active inequality constraints.

Then
\[
x^*
\]
is a local maximizer of
\[
f
\]
subject to
\[
g(x)=0
\]
and
\[
h(x)\le0.
\]

\paragraph{Proof.}

Convert the inequality-constrained problem into an equality-constrained problem by replacing each inequality constraint
\[
h_k(x)\le0
\]
with
\[
h_k(x)+y_k^2=0,
\qquad
k=1,\dots,K.
\]

Then apply the proof used for the second-order sufficiency conditions of the equality-constrained problem.

\hfill\(\blacksquare\)

\section{Parametric Estimation}

\subsection{OLS Estimation}

We may analyze a production system by estimating a single cost or profit function.

Specify the model as
\[
y=X\beta+u,
\]
where:

\begin{itemize}
    \item \(y\) is a \(T\times 1\) vector of observed dependent variables, for example \(c(w,y)\);
    \item \(X\) is a \(T\times K\) matrix of observed explanatory variables, including a constant term;
    \item \(\beta\) is a \(K\times 1\) vector of parameters to be estimated;
    \item \(u\) is a \(T\times 1\) vector of disturbances.
\end{itemize}

We impose the following assumptions of the classical linear model.

\[
E(u_t\mid X)=0
\]
for all \(t\). This is strict exogeneity.

\[
\operatorname{rank}(X)=K.
\]
This means that there is no perfect multicollinearity.

\[
E(u_t^2\mid X)=\sigma^2.
\]
This is homoskedasticity.

\[
E(u_tu_s\mid X)=0,
\qquad
t\ne s.
\]
This means that there is no correlation between observations.

OLS minimizes the sum of squared residuals:
\[
\hat\beta
=
\arg\min_{\tilde\beta}
SSR(\tilde\beta),
\]
where
\[
SSR(\tilde\beta)
=
\sum_{t=1}^T
(y_t-x_t^\top\tilde\beta)^2.
\]

In matrix form,
\[
SSR(\tilde\beta)
=
(y-X\tilde\beta)^\top
(y-X\tilde\beta).
\]

Expanding,
\[
SSR(\tilde\beta)
=
y^\top y
-
2y^\top X\tilde\beta
+
\tilde\beta^\top X^\top X\tilde\beta.
\]

The first-order condition is
\[
\frac{\partial SSR(\tilde\beta)}{\partial \tilde\beta}
=
-2X^\top y
+
2X^\top X\tilde\beta
=
0.
\]

Thus,
\[
X^\top X\hat\beta
=
X^\top y.
\]

Since
\[
\operatorname{rank}(X)=K,
\]
the matrix
\[
X^\top X
\]
is invertible.

Therefore,
\[
\hat\beta
=
(X^\top X)^{-1}X^\top y.
\]



Thus, the OLS estimator is
\[
\hat\beta
=
(X^\top X)^{-1}X^\top y.
\]

The OLS estimator of \(\sigma^2\) is
\[
s^2
=
\frac{\hat u^\top \hat u}{T-K}
=
\frac{SSR}{T-K}.
\]

The OLS estimator is unbiased:
\[
E(\hat\beta\mid X)=\beta.
\]

It is also efficient within the class of linear unbiased estimators.

For hypothesis testing, we need an additional assumption:
\[
u\mid X\sim N(0,\sigma^2I_T).
\]

\subsubsection{Hypothesis Testing for Individual Coefficients}

Consider the null hypothesis
\[
H_0:\beta_k=\bar\beta_k.
\]

The test statistic is
\[
t_k
=
\frac{\hat\beta_k-\bar\beta_k}{SE(\hat\beta_k)}.
\]

Since
\[
\operatorname{Var}(\hat\beta\mid X)
=
\sigma^2(X^\top X)^{-1},
\]
we estimate
\[
SE(\hat\beta_k)
=
\sqrt{
s^2
\left[
(X^\top X)^{-1}
\right]_{kk}
}.
\]

Thus,
\[
t_k
=
\frac{
\hat\beta_k-\bar\beta_k
}{
\sqrt{
s^2
\left[
(X^\top X)^{-1}
\right]_{kk}
}
}.
\]

Under the null hypothesis,
\[
t_k\sim t_{T-K}.
\]

\subsubsection{Testing a Set of Linear Hypotheses}

Consider
\[
H_0:R\beta=r,
\]
where \(R\) is a \(J\times K\) matrix and \(r\) is a \(J\times 1\) vector.

For example, suppose
\[
K=4,
\]
and the restrictions are
\[
\beta_2=\beta_3,
\qquad
\beta_4=0.
\]

Then
\[
R
=
\begin{bmatrix}
0 & 1 & -1 & 0 \\
0 & 0 & 0 & 1
\end{bmatrix},
\qquad
r
=
\begin{bmatrix}
0\\
0
\end{bmatrix}.
\]

The Wald test statistic is
\[
F
=
\frac{
(R\hat\beta-r)^\top
\left[
R(X^\top X)^{-1}R^\top
\right]^{-1}
(R\hat\beta-r)/J
}{
s^2
}.
\]

Under the null hypothesis,
\[
F\sim F(J,T-K).
\]

Equivalently, using the likelihood-ratio principle, we can write
\[
F
=
\frac{
(SSR_R-SSR_U)/J
}{
SSR_U/(T-K)
},
\]
where:

\begin{itemize}
    \item \(J\) is the number of restrictions,
    \item \(SSR_R\) is the sum of squared residuals from the restricted model,
    \item \(SSR_U\) is the sum of squared residuals from the unrestricted model.
\end{itemize}

\subsubsection{Nerlove's Cost Function Estimation}

Nerlove (1963) is a path-breaking work in applied production analysis.

The paper showed that one can recover all information on the production function by estimating a cost function.

Let:

\begin{itemize}
    \item \(y\): electricity output, measured in kwh;
    \item \(x_1\): labor input;
    \item \(x_2\): capital input;
    \item \(x_3\): fuel input;
    \item \(p_1\): wage rate;
    \item \(p_2\): capital price;
    \item \(p_3\): fuel price;
    \item \(u\): residual.
\end{itemize}

Nerlove specified the Cobb--Douglas production function:
\[
y
=
a_0x_1^{\alpha_1}
x_2^{\alpha_2}
x_3^{\alpha_3}
u.
\]

From this Cobb--Douglas production function, the corresponding cost function can be derived as
\[
c
=
ky^{1/r}
p_1^{\alpha_1/r}
p_2^{\alpha_2/r}
p_3^{\alpha_3/r}
v,
\]
where
\[
r=\alpha_1+\alpha_2+\alpha_3,
\]
\[
v=u^{-1/r},
\]
and
\[
k
=
r
\left(
a_0
\alpha_1^{\alpha_1}
\alpha_2^{\alpha_2}
\alpha_3^{\alpha_3}
\right)^{-1/r}.
\]

Taking logarithms gives
\[
\ln c
=
K
+
\frac{1}{r}\ln y
+
\frac{\alpha_1}{r}\ln p_1
+
\frac{\alpha_2}{r}\ln p_2
+
\frac{\alpha_3}{r}\ln p_3
+
v.
\]

Before estimation, we impose the homogeneity condition.

This can be done by transforming variables:
\[
\ln c'
=
\ln c-\ln p_3,
\]
\[
\ln p_1'
=
\ln p_1-\ln p_3,
\]
and
\[
\ln p_2'
=
\ln p_2-\ln p_3.
\]

That is, we estimate
\[
\ln c'
=
K
+
\frac{1}{r}\ln y
+
\frac{\alpha_1}{r}\ln p_1'
+
\frac{\alpha_2}{r}\ln p_2'
+
v.
\]

The following is a STATA program for estimating the model:
\[
\begin{aligned}
&\texttt{\#delimit;}\\
&\texttt{use "C:\textbackslash WORK\textbackslash class\textbackslash micro\textbackslash nerlove.dta", clear;}\\
&\texttt{gen ltc = ln(tc);}\\
&\texttt{gen lq = ln(q);}\\
&\texttt{gen lp1 = ln(p1);}
\end{aligned}
\]


\[
\texttt{gen lpf = ln(pf);}
\]
\[
\texttt{gen lpk = ln(pk);}
\]
\[
\texttt{gen ltc2 = ltc - lpf;}
\]
\[
\texttt{gen lpl2 = lpl - lpf;}
\]
\[
\texttt{gen lpk2 = lpk - lpf;}
\]
\[
\texttt{regress ltc2 lpl2 lpk2 lq;}
\]

The estimated equation is
\[
\ln c-\ln p_3
=
-4.7
+
0.59(\ln p_1-\ln p_3)
-
0.007(\ln p_2-\ln p_3)
+
0.72\ln y.
\]

The standard errors are
\[
(0.885)
\qquad
(0.205)
\qquad
(0.191)
\qquad
(0.017).
\]

The coefficient of determination is
\[
R^2=0.932.
\]

The following command conducts an \(F\)-test for the constant returns to scale hypothesis:
\[
H_0:r=1.
\]

Since the coefficient on \(\ln y\) is
\[
\frac{1}{r},
\]
the CRS hypothesis implies
\[
\frac{1}{r}=1.
\]

Thus, the STATA command is
\[
\texttt{test lq = 1;}
\]

The test statistic is
\[
F(1,141)=256.63.
\]

Therefore,
\[
H_0
\]
is rejected.

\subsection{SUR Estimation}

Estimating the entire cost function using OLS requires a large number of observations when we specify a flexible functional form.

Instead, we can estimate a system of shorter equations by applying Shephard's lemma to the cost function.

Suppose we want to estimate the following translog cost function:
\[
\begin{aligned}
\ln c
&=
\beta_0
+
\sum_{i=1}^3 \beta_i\ln p_i
+
\frac{1}{2}
\sum_{i=1}^3
\sum_{j=1}^3
\beta_{ij}\ln p_i\ln p_j \\
&\quad
+
\beta_y\ln y
+
\frac{1}{2}\beta_{yy}(\ln y)^2
+
\sum_{i=1}^3
\beta_{iy}\ln p_i\ln y
+
u.
\end{aligned}
\]

The parameter restrictions are
\[
\beta_{ij}=\beta_{ji},
\]
\[
\sum_{i=1}^3\beta_i=1,
\]
\[
\sum_{i=1}^3\beta_{ij}=0,
\qquad
j=1,2,3,
\]
\[
\sum_{j=1}^3\beta_{ij}=0,
\qquad
i=1,2,3,
\]
and
\[
\sum_{i=1}^3\beta_{iy}=0.
\]

By Shephard's lemma,
\[
s_i
=
\frac{p_i x_i}{c}
=
\frac{\partial \ln c}{\partial \ln p_i}.
\]

Therefore, the cost-share equations are
\[
s_1
=
\beta_1
+
\beta_{11}\ln p_1
+
\beta_{12}\ln p_2
+
\beta_{13}\ln p_3
+
\beta_{1y}\ln y,
\]
\[
s_2
=
\beta_2
+
\beta_{12}\ln p_1
+
\beta_{22}\ln p_2
+
\beta_{23}\ln p_3
+
\beta_{2y}\ln y,
\]
and
\[
s_3
=
\beta_3
+
\beta_{13}\ln p_1
+
\beta_{23}\ln p_2
+
\beta_{33}\ln p_3
+
\beta_{3y}\ln y.
\]

Instead of estimating
\[
\ln c
\]
directly, we may estimate the share equations.

Alternatively, we may estimate the share equations and the equation for
\[
\ln c
\]
simultaneously.

Since
\[
s_1+s_2+s_3=1
\]
identically, we usually drop one of the share equations.

Suppose we have \(N\) equations.

Represent the \(i\)-th equation as
\[
y_i=X_i\beta_i+u_i,
\qquad
i=1,\dots,N.
\]

We impose the following assumptions.

\[
E(u_{it}\mid X_t)=0.
\]

\[
\operatorname{Var}(u_{it}\mid X_t)
=
E(u_{it}^2\mid X_t)
=
\sigma_i^2.
\]

This variance is constant over observations but may differ across equations.

For different equations,
\[
\operatorname{Cov}(u_{it},u_{jt}\mid X_t)
=
E(u_{it}u_{jt}\mid X_t)
=
\sigma_{ij}.
\]


This is called contemporaneous correlation.

We also assume
\[
\operatorname{Cov}(u_{it},u_{js}\mid X)
=
E(u_{it}u_{js}\mid X)
=
0,
\qquad
t\ne s.
\]

This means that there is no autocorrelation across observations.

Equivalently,
\[
E(u_i\mid X)=0,
\]
and
\[
E(u_iu_j^\top\mid X)=\sigma_{ij}I_T.
\]

Now combine all equations into one system:
\[
y=X\beta+u,
\]
where
\[
y=
\begin{bmatrix}
y_1\\
y_2\\
\vdots\\
y_N
\end{bmatrix},
\qquad
X=
\begin{bmatrix}
X_1 & 0 & \cdots & 0\\
0 & X_2 & \cdots & 0\\
\vdots & \vdots & \ddots & \vdots\\
0 & 0 & \cdots & X_N
\end{bmatrix},
\]
\[
\beta=
\begin{bmatrix}
\beta_1\\
\beta_2\\
\vdots\\
\beta_N
\end{bmatrix},
\qquad
u=
\begin{bmatrix}
u_1\\
u_2\\
\vdots\\
u_N
\end{bmatrix}.
\]

The dimensions are
\[
y:NT\times 1,
\qquad
X:NT\times K,
\qquad
\beta:K\times 1,
\qquad
u:NT\times 1,
\]
where
\[
K=\sum_{i=1}^N K_i.
\]

The covariance matrix for the complete disturbance vector is
\[
\Phi
=
E(uu^\top)
=
\begin{bmatrix}
\sigma_{11}I_T & \sigma_{12}I_T & \cdots & \sigma_{1N}I_T\\
\sigma_{21}I_T & \sigma_{22}I_T & \cdots & \sigma_{2N}I_T\\
\vdots & \vdots & \ddots & \vdots\\
\sigma_{N1}I_T & \sigma_{N2}I_T & \cdots & \sigma_{NN}I_T
\end{bmatrix}.
\]

Equivalently,
\[
\Phi
=
\Sigma\otimes I_T,
\]
where
\[
\Sigma
=
\begin{bmatrix}
\sigma_{11} & \sigma_{12} & \cdots & \sigma_{1N}\\
\sigma_{21} & \sigma_{22} & \cdots & \sigma_{2N}\\
\vdots & \vdots & \ddots & \vdots\\
\sigma_{N1} & \sigma_{N2} & \cdots & \sigma_{NN}
\end{bmatrix},
\]
with
\[
\sigma_{ij}=\sigma_{ji}.
\]

Assume that
\[
\Sigma
\]
and
\[
\Phi
\]
are positive definite.

Therefore, there exists a matrix \(P\) such that
\[
P\Phi P^\top=I_{NT}.
\]

Using this matrix, we transform the model:
\[
Py=PX\beta+Pu.
\]

Define
\[
y^*=Py,
\qquad
X^*=PX,
\qquad
u^*=Pu.
\]

Then
\[
y^*=X^*\beta+u^*.
\]

The transformed disturbance vector has mean
\[
E(u^*)=E(Pu)=PE(u)=0.
\]

Its covariance matrix is
\[
E(u^*u^{*\top})
=
E(Puu^\top P^\top)
=
PE(uu^\top)P^\top
=
P\Phi P^\top
=
I_{NT}.
\]

Thus, the elements of the transformed disturbance vector are uncorrelated and have identical variances.

The least squares estimator in the transformed model is
\[
\hat\beta
=
(X^{*\top}X^*)^{-1}X^{*\top}y^*.
\]

Writing this estimator in terms of the original observations gives
\[
\hat\beta
=
(X^\top P^\top PX)^{-1}X^\top P^\top Py.
\]

Since
\[
P\Phi P^\top=I_{NT},
\]
we have
\[
\Phi=P^{-1}(P^\top)^{-1}.
\]

Hence,
\[
\Phi^{-1}=P^\top P.
\]

Therefore,
\[
\hat\beta
=
(X^\top\Phi^{-1}X)^{-1}
X^\top\Phi^{-1}y.
\]

Since
\[
\Phi^{-1}
=
\Sigma^{-1}\otimes I_T,
\]
we can also write
\[
\hat\beta
=
\left[
X^\top(\Sigma^{-1}\otimes I_T)X
\right]^{-1}
X^\top(\Sigma^{-1}\otimes I_T)y.
\]

This is the best linear unbiased estimator when \(\Sigma\) is known.

In practice, however, \(\Sigma\) is unknown and must be estimated.

To estimate \(\sigma_{ij}\), first estimate each equation by ordinary least squares:
\[
\hat\beta_i
=
(X_i^\top X_i)^{-1}X_i^\top y_i.
\]

Then obtain the least squares residuals:
\[
\hat e_i
=
y_i-X_i\hat\beta_i.
\]

A consistent estimator of the covariance between equations \(i\) and \(j\) is
\[
\hat\sigma_{ij}
=
\frac{1}{T}
\hat e_i^\top \hat e_j.
\]

Equivalently,
\[
\hat\sigma_{ij}
=
\frac{1}{T}
\sum_{t=1}^T
\hat e_{it}\hat e_{jt}.
\]

Define
\[
\hat\Sigma
\]
as the matrix obtained from
\[
\Sigma
\]
by replacing each
\[
\sigma_{ij}
\]
with
\[
\hat\sigma_{ij}.
\]

Then Zellner's seemingly unrelated regression estimator is
\[
\hat\beta_{SUR}
=
\left[
X^\top(\hat\Sigma^{-1}\otimes I_T)X
\right]^{-1}
X^\top(\hat\Sigma^{-1}\otimes I_T)y.
\]

The SUR estimator can therefore be written as
\[
\hat\beta
=
\left[
X^\top
(\hat\Sigma^{-1}\otimes I_T)
X
\right]^{-1}
X^\top
(\hat\Sigma^{-1}\otimes I_T)
y.
\]

Alternatively, we may estimate
\[
\sigma_{ij}
\]
and
\[
\beta
\]
iteratively.

Starting from an initial estimator,
new variance estimates are computed as
\[
\hat\sigma_{ij}
=
\frac{1}{T}
(y_i-X_i\hat\beta_i)^\top
(y_j-X_j\hat\beta_j).
\]

Equivalently,
\[
\hat\sigma_{ij}
=
\frac{1}{T}
\sum_{t=1}^T
\hat e_{it}\hat e_{jt}.
\]

Using these updated covariance estimates,
we form a new estimator of
\[
\beta.
\]

The process is repeated until convergence.

When the disturbances follow a multivariate normal distribution,
this iterative estimator coincides with the maximum likelihood estimator.

\subsection{Hypothesis Testing in SUR}

Suppose we wish to test the linear hypotheses
\[
H_0:R\beta=r.
\]

Unlike the classical OLS case,
the exact finite-sample distribution of the test statistic is not available because the covariance matrix
\[
\Sigma
\]
must itself be estimated.

However, asymptotically,
the following quadratic form has a chi-square distribution under
\[
H_0:
\]
\[
g
=
(R\hat\beta-r)^\top
\left[
R
\left(
X^\top
(\hat\Sigma^{-1}\otimes I_T)
X
\right)^{-1}
R^\top
\right]^{-1}
(R\hat\beta-r).
\]

Under the null hypothesis,
\[
g
\overset{a}{\sim}
\chi^2_J,
\]
where
\[
J
\]
is the number of restrictions.

Alternatively,
we may use an asymptotic
\[
F
\]
test.

Since
\[
JF(J,NT-K)
\overset{a}{\sim}
\chi^2_J,
\]
the statistic
\[
AF
=
\frac{g}{J}
\]
has approximately an
\[
F(J,NT-K)
\]
distribution.

\subsection{Christensen and Greene (1976)}

Christensen and Greene (1976, \emph{Journal of Political Economy})
extended the work of Nerlove (1963).

They estimated a system of conditional factor demand equations derived from a translog cost function.

Consider the translog cost function:
\[
\ln c
=
\beta_0
+
\sum_{i=1}^3
\beta_i\ln p_i
+
\frac12
\sum_{i=1}^3
\sum_{j=1}^3
\beta_{ij}\ln p_i\ln p_j
+
\beta_y\ln y
+
\frac12
\beta_{yy}(\ln y)^2
+
\sum_{i=1}^3
\beta_{iy}\ln p_i\ln y
+
u.
\]

The parameter restrictions are
\[
\beta_{ij}=\beta_{ji},
\]
\[
\sum_{i=1}^3\beta_i=1,
\]
\[
\sum_{i=1}^3\beta_{ij}=0,
\qquad
j=1,2,3,
\]
and
\[
\sum_{i=1}^3\beta_{iy}=0.
\]

Applying Shephard's lemma,
the factor share equations become
\[
s_i
=
\frac{p_ix_i}{c}
=
\beta_i
+
\sum_{j=1}^3
\beta_{ij}\ln p_j
+
\beta_{iy}\ln y,
\qquad
i=1,2,3.
\]

Explicitly,
\[
s_1
=
\beta_1
+
\beta_{11}\ln p_1
+
\beta_{12}\ln p_2
+
\beta_{13}\ln p_3
+
\beta_{1y}\ln y,
\]
\[
s_2
=
\beta_2
+
\beta_{12}\ln p_1
+
\beta_{22}\ln p_2
+
\beta_{23}\ln p_3
+
\beta_{2y}\ln y,
\]
\[
s_3
=
\beta_3
+
\beta_{13}\ln p_1
+
\beta_{23}\ln p_2
+
\beta_{33}\ln p_3
+
\beta_{3y}\ln y.
\]

Since the adding-up condition implies
\[
s_1+s_2+s_3=1,
\]
the covariance matrix of the disturbances is singular if all three equations are estimated simultaneously.

Therefore,
one share equation is usually dropped.

For example,
we may estimate only the labor and capital share equations while omitting the fuel share equation.

The remaining equations can then be estimated jointly using SUR.

This approach substantially reduces the number of parameters relative to direct estimation of the full cost function,
while still preserving the theoretical restrictions implied by duality theory.


\section{Proof of Theorem 3}

We derive the Linear Expenditure System (LES) from the conditions of homogeneity, symmetry, adding-up, and negativity.

\subsection{Homogeneity}

Suppose the demand system takes the form
\[
x_i
=
\sum_{j=1}^N
a_{ij}\frac{p_j}{p_i}
+
\beta_i\frac{m}{p_i}.
\tag{1}
\]

Because demand functions are homogeneous of degree zero in prices and income,
dividing all prices and income by
\[
p_i
\]
does not change
\[
x_i.
\]

Hence the above representation is consistent with homogeneity.

\subsection{Symmetry}

Symmetry of the Slutsky matrix requires
\[
\frac{\partial x_i}{\partial p_j}
+
x_j\frac{\partial x_i}{\partial m}
=
\frac{\partial x_j}{\partial p_i}
+
x_i\frac{\partial x_j}{\partial m}.
\]

From equation (1),
\[
\frac{\partial x_i}{\partial p_j}
=
\frac{a_{ij}}{p_i},
\qquad
\frac{\partial x_i}{\partial m}
=
\frac{\beta_i}{p_i}.
\]

Similarly,
\[
\frac{\partial x_j}{\partial p_i}
=
\frac{a_{ji}}{p_j},
\qquad
\frac{\partial x_j}{\partial m}
=
\frac{\beta_j}{p_j}.
\]

Substituting into the symmetry condition gives
\[
\frac{a_{ij}}{p_i}
+
x_j\frac{\beta_i}{p_i}
=
\frac{a_{ji}}{p_j}
+
x_i\frac{\beta_j}{p_j}.
\]

Multiplying both sides by
\[
p_ip_j
\]
yields
\[
a_{ij}p_j+\beta_ip_jx_j
=
a_{ji}p_i+\beta_jp_ix_i.
\tag{2}
\]

Using equation (1),
\[
p_ix_i
=
\sum_{k=1}^N a_{ik}p_k+\beta_im,
\]
and similarly,
\[
p_jx_j
=
\sum_{k=1}^N a_{jk}p_k+\beta_jm.
\]

Substituting these into equation (2),
\[
a_{ij}p_j
+
\beta_i
\left(
\sum_{k=1}^N a_{jk}p_k+\beta_jm
\right)
=
a_{ji}p_i
+
\beta_j
\left(
\sum_{k=1}^N a_{ik}p_k+\beta_im
\right).
\]

The income terms cancel:
\[
\beta_i\beta_jm-\beta_j\beta_im=0.
\]

Therefore,
\[
a_{ij}p_j
+
\beta_i\sum_{k=1}^N a_{jk}p_k
=
a_{ji}p_i
+
\beta_j\sum_{k=1}^N a_{ik}p_k.
\]

Rearranging terms,
\[
(a_{ij}+\beta_ia_{jj})p_j
+
\beta_ia_{ji}p_i
+
\beta_i\sum_{k\ne i,j}a_{jk}p_k
\]
\[
=
(a_{ji}+\beta_ja_{ii})p_i
+
\beta_ja_{ij}p_j
+
\beta_j\sum_{k\ne i,j}a_{ik}p_k.
\]

Since this equality must hold for all price vectors,
the coefficients on corresponding prices must be identical.

Thus,
\[
a_{ij}
=
\beta_i
\frac{a_{jj}}{\beta_j-1},
\]
\[
a_{ii}
=
\beta_i
\frac{a_{ji}}{\beta_j}
-
\frac{a_{ij}}{\beta_j},
\]
and
\[
a_{ik}
=
\beta_i
\frac{a_{jk}}{\beta_j},
\qquad
i\ne j\ne k.
\]

Define
\[
\gamma_k
=
\frac{a_{jk}}{\beta_j},
\qquad
k\ne j,
\]
and
\[
\gamma_j
=
\frac{a_{jj}}{\beta_j-1}.
\]

Then the general form becomes
\[
a_{ik}
=
\beta_i
(\gamma_k-b_{ik}\gamma_i),
\]
where
\[
b_{ik}
=
\begin{cases}
1, & i=k,\\
0, & i\ne k.
\end{cases}
\]

\subsection{Adding-Up}

From equation (1),
\[
p_ix_i
=
\sum_{j=1}^N a_{ij}p_j+\beta_im.
\]

Substituting the symmetry representation,
\[
p_ix_i
=
\sum_{j=1}^N
\left[
\beta_i(\gamma_j-b_{ij}\gamma_i)
\right]p_j
+
\beta_im.
\]

Thus,
\[
p_ix_i
=
\beta_i\sum_{j=1}^N p_j\gamma_j
-
\beta_ip_i\gamma_i
+
\beta_im.
\tag{3}
\]

Adding equation (3) across all goods,
\[
m
=
\sum_{i=1}^N p_ix_i.
\]

Hence,
\[
m
=
\sum_{i=1}^N
\beta_i
\sum_{j=1}^N p_j\gamma_j
-
\sum_{i=1}^N \beta_ip_i\gamma_i
+
m\sum_{i=1}^N\beta_i.
\]

Rearranging,
\[
\sum_{j=1}^N p_j\gamma_j
\left(
\sum_{i=1}^N\beta_i-1
\right)
+
m
\left(
\sum_{i=1}^N\beta_i-1
\right)
=
0.
\]

Since this must hold for arbitrary prices and income,
we require
\[
\sum_{i=1}^N\beta_i=1.
\]

Using this condition,
equation (3) becomes
\[
p_ix_i
=
-\beta_ip_i\gamma_i
+
\beta_i
\left(
m+\sum_{j=1}^N p_j\gamma_j
\right).
\]

Redefining constants,
\[
x_i
=
\gamma_i
+
\frac{\beta_i}{p_i}
\left(
m-\sum_{j=1}^N p_j\gamma_j
\right),
\]
with
\[
\sum_{i=1}^N\beta_i=1.
\]

Therefore,
the demand system becomes
\[
p_ix_i
=
p_i\gamma_i
+
\beta_i
\left(
m-\sum_{j=1}^N p_j\gamma_j
\right).
\]

This is the Linear Expenditure System (LES).

\subsection{Negativity}

Finally,
the Slutsky substitution matrix must be negative semidefinite.

The own-price derivative is
\[
\frac{\partial x_i}{\partial p_i}
+
x_i\frac{\partial x_i}{\partial m}.
\]

Using the LES form,
\[
\frac{\partial x_i}{\partial p_i}
=
-
\frac{\beta_i}{p_i^2}
\left(
m-\sum_{j=1}^N p_j\gamma_j
\right)
-
\frac{\beta_i\gamma_i}{p_i},
\]
and
\[
\frac{\partial x_i}{\partial m}
=
\frac{\beta_i}{p_i}.
\]

Therefore,
\[
\frac{\partial x_i}{\partial p_i}
+
x_i\frac{\partial x_i}{\partial m}
=
-
\frac{\beta_i}{p_i}(x_i-\gamma_i)
+
\frac{\beta_i}{p_i}x_i.
\]

Simplifying,
\[
=
\frac{\beta_i\gamma_i}{p_i}
-
\frac{x_i}{p_i}.
\]

Hence negativity requires
\[
x_i>\gamma_i,
\]
and
\[
0<\beta_i<1.
\]

Together with the adding-up condition,
the final restrictions are
\[
x_i>\gamma_i,
\qquad
0<\beta_i<1,
\qquad
\sum_{i=1}^N\beta_i=1.
\]

These are the conditions characterizing the Linear Expenditure System.



\section{Estimation of a System of Nonlinear Equations}

Many applied demand systems are nonlinear in parameters.

As a result,
we often need to estimate a system of nonlinear equations simultaneously.

Suppose we wish to estimate the following system:
\[
Y_1=f_1(X,\beta)+e_1,
\]
\[
Y_2=f_2(X,\beta)+e_2,
\]
\[
\vdots
\]
\[
Y_N=f_N(X,\beta)+e_N.
\]

Here:

\begin{itemize}
    \item \(Y_i\) is the dependent variable in equation \(i\),
    \item \(f_i(\cdot)\) is a nonlinear function,
    \item \(X\) denotes the observed explanatory variables,
    \item \(\beta\) is the vector of parameters,
    \item \(e_i\) is the disturbance term.
\end{itemize}

Define the disturbance vector
\[
e^\top
=
(e_1^\top,\dots,e_N^\top).
\]

Assume
\[
E(ee^\top)
=
\Sigma\otimes I_T,
\]
where
\[
\Sigma
=
\begin{bmatrix}
\sigma_{11} & \sigma_{12} & \cdots & \sigma_{1N}\\
\sigma_{21} & \sigma_{22} & \cdots & \sigma_{2N}\\
\vdots & \vdots & \ddots & \vdots\\
\sigma_{N1} & \sigma_{N2} & \cdots & \sigma_{NN}
\end{bmatrix}
\]
is the covariance matrix across equations.

The
\[
(i,j)
\]
element of
\[
\Sigma
\]
is
\[
\sigma_{ij},
\]
and
\[
E(e_ie_j^\top)=\sigma_{ij}I_T.
\]

Each equation may have a different nonlinear functional form.

\subsection{Maximum Likelihood Estimation}

One common estimation method is maximum likelihood estimation (MLE).

Assume additionally that the disturbances are jointly normally distributed.

Then the log-likelihood function for
\[
\beta
\]
and
\[
\Sigma
\]
is
\[
\ln L(\beta,\Sigma)
=
-\frac{NT}{2}\ln(2\pi)
-
\frac12\ln|\Sigma\otimes I_T|
-
\frac12
e^\top
(\Sigma^{-1}\otimes I_T)
e.
\]

Using properties of determinants,
\[
|\Sigma\otimes I_T|
=
|\Sigma|^T.
\]

Therefore,
\[
\ln|\Sigma\otimes I_T|
=
T\ln|\Sigma|.
\]

Hence the likelihood becomes
\[
\ln L(\beta,\Sigma)
=
-\frac{NT}{2}\ln(2\pi)
-
\frac{T}{2}\ln|\Sigma|
-
\frac12
e^\top
(\Sigma^{-1}\otimes I_T)
e.
\]

Using the trace operator,
the quadratic form can be rewritten as
\[
e^\top
(\Sigma^{-1}\otimes I_T)
e
=
\operatorname{tr}(S\Sigma^{-1}),
\]
where
\[
S
=
\begin{bmatrix}
e_1^\top e_1 & \cdots & e_1^\top e_N\\
\vdots & \ddots & \vdots\\
e_N^\top e_1 & \cdots & e_N^\top e_N
\end{bmatrix}.
\]

The
\[
(i,j)
\]
element of
\[
S
\]
is
\[
e_i^\top e_j
=
[Y_i-f_i(X,\beta)]^\top
[Y_j-f_j(X,\beta)].
\]

Thus,
the log-likelihood function becomes
\[
\ln L(\beta,\Sigma)
=
-\frac{NT}{2}\ln(2\pi)
-
\frac{T}{2}\ln|\Sigma|
-
\frac12
\operatorname{tr}(S\Sigma^{-1}).
\]

The MLE chooses
\[
\hat\beta
\]
and
\[
\hat\Sigma
\]
to maximize this expression.

\subsection{Estimator of \(\Sigma\)}

Differentiate the log-likelihood function with respect to
\[
\Sigma^{-1}.
\]

Using matrix derivative results,
\[
\frac{\partial \ln|\Sigma|}{\partial \Sigma^{-1}}
=
-\Sigma,
\]
and
\[
\frac{\partial}{\partial \Sigma^{-1}}
\operatorname{tr}(S\Sigma^{-1})
=
S.
\]

The first-order condition is therefore
\[
\frac{T}{2}\Sigma
-
\frac12 S
=
0.
\]

Hence,
\[
\hat\Sigma
=
\frac{S}{T}.
\]

Substituting this estimator into the log-likelihood function yields the concentrated log-likelihood function:
\[
\ln L^*(\beta)
=
\text{constant}
-
\frac{T}{2}\ln|S|.
\]

Therefore,
the MLE of
\[
\beta
\]
is the value minimizing
\[
|S|.
\]

Explicitly,
\[
|S|
=
\begin{vmatrix}
e_1^\top e_1 & \cdots & e_1^\top e_N\\
\vdots & \ddots & \vdots\\
e_N^\top e_1 & \cdots & e_N^\top e_N
\end{vmatrix}.
\]

\subsection{Covariance Matrix of the MLE}

The asymptotic covariance matrix of
\[
\hat\beta
\]
is
\[
V_{\hat\beta}
=
\left[
\sum_{t=1}^T
\left(
\frac{\partial f_t}{\partial\beta}
\right)^\top
(\Sigma^{-1}\otimes I_T)
\left(
\frac{\partial f_t}{\partial\beta}
\right)
\right]^{-1},
\]
where
\[
\frac{\partial f_t}{\partial\beta}
\]
is the Jacobian matrix of the nonlinear system evaluated at observation
\[
t.
\]

\subsection{Exercise}

Prove the following two identities:

\[
\frac12\ln|\Sigma\otimes I_T|
=
\frac{T}{2}\ln|\Sigma|,
\]
and
\[
\frac12
e^\top
(\Sigma^{-1}\otimes I_T)
e
=
\frac12
\operatorname{tr}(S\Sigma^{-1}).
\]

\section{Examples}

Suppose we observe household expenditure data across seven occupational categories over the years 1982--2003.

The ten expenditure categories are:

\begin{enumerate}
    \item Food
    \item Housing
    \item Heating and water
    \item Furniture and housekeeping
    \item Clothing
    \item Health and medical services
    \item Education
    \item Entertainment
    \item Transportation and telecommunication
    \item Miscellaneous expenditures
\end{enumerate}

These expenditure shares may be modeled jointly using nonlinear demand systems such as:

\begin{itemize}
    \item LES
    \item AIDS
    \item Translog demand systems
    \item QUAIDS
\end{itemize}

Because the share equations are contemporaneously correlated,
joint estimation methods such as nonlinear SUR or full-information maximum likelihood are generally more efficient than estimating each equation separately.


